% Glossaire des modèles (type=modeles) : détecteurs, backbones, autoencodeurs,
% apprentissage continu, opérateurs et réseaux quantifiés.
% Chiffres : \citep[§]{paper_id} s'ils viennent d'une fiche ou de benchmarks.csv,
% sinon \horscorpus. Les parents sont définis avant leurs enfants.

% ---------------------------------------------------------------- Famille YOLO
\newglossaryentry{yolo}{type=modeles, name={YOLO}, sort={YOLO},
  description={\emph{You Only Look Once} : famille de détecteurs en une étape, qui prédisent
  en une seule passe, sur une grille, des boîtes, un score d'objet et des scores de classe,
  puis filtrent les doublons par NMS. Depuis YOLOv3, l'architecture se découpe en backbone
  (extraction de caractéristiques), neck (fusion multi-échelle) et tête de détection. C'est
  le détecteur dominant du corpus : de Tiny YOLO sur Virtex-7 (REQ-YOLO) à YOLOv26n sur DPU,
  et le cas d'application visé par la thèse (détection de personnes).}}

\newglossaryentry{yolov2}{type=modeles, name={YOLOv2}, sort={YOLO 02}, parent=yolo,
  description={Deuxième version (YOLO9000) : backbone Darknet-19, boîtes d'ancrage
  (\emph{anchors}) obtenues par k-means, normalisation par lot, couche \emph{passthrough}
  (reorg) qui ramène des caractéristiques fines, une seule échelle de sortie ; entrée usuelle
  416$\times$416. Face à YOLOv1, elle ajoute les ancres et supprime les couches denses. Dans le
  corpus, cible d'accélérateurs HLS sur Zynq : Zhang 2024 (YOLOv2 et v2-tiny) et Kim 2025, dont
  le YOLOv2 compte 29,5\,GFLOPs mais ne tient que 2,33\,GOPS sur XC7Z020
  \citep[Table~9]{2025-kim-design-implementation-yolov2-accelerator-zynq}.}}

\newglossaryentry{yolov2-tiny}{type=modeles, name={YOLOv2-tiny}, sort={YOLO 02 tiny},
  parent=yolov2,
  description={Version réduite de YOLOv2 (\enquote{Tiny YOLO}) : pile de convolutions
  3$\times$3 et de max-pooling sans Darknet-19, sans passthrough, une seule échelle. REQ-YOLO
  la compresse (9 couches de convolution, entrée 416$\times$416)
  \citep[§2.1]{2019-ding-req-yolo-resource-aware-efficient} ; FINN-R en dérive Tincy YOLO, et
  CoDeNet la prend comme référence (57,1 d'AP50 VOC en FP32)
  \citep[Table~3]{arxiv-2006.08357-codenet-efficient-deployment-input-adaptive}.}}

\newglossaryentry{yolov3}{type=modeles, name={YOLOv3}, sort={YOLO 03}, parent=yolo,
  description={Troisième version : backbone Darknet-53 (blocs résiduels), prédiction à trois
  échelles par un neck de type FPN, trois ancres par échelle, classes en sigmoïdes
  indépendantes. Face à YOLOv2 : multi-échelle et backbone plus profond, d'où de meilleurs
  petits objets. Dans le corpus : version élaguée sur Zynq XC7Z035 avec suivi Deepsort (Zhai
  2023), tête de MCUNetV2 sur MCU, base de Gaussian YOLOv3, mesures sur GPU embarqué (Feng
  2022).}}

\newglossaryentry{yolov3-tiny}{type=modeles, name={YOLOv3-tiny}, sort={YOLO 03 tiny},
  parent=yolov3,
  description={Version réduite de YOLOv3 : backbone de convolutions 3$\times$3 et de
  max-pooling, deux échelles de sortie au lieu de trois. Détecteur le plus porté sur FPGA dans
  le corpus : SATAY, Zhang 2024, Zhai 2023, la référence w4a4 de LP-YOLO (0,761 de mAP@0,5 VOC
  à 12,35\,FPS sur PYNQ-Z1 \citep[§V]{arxiv-2609.24757-npu-accelerator-quantized-real-time}\pp),
  Fata 2026 (élagué à 70\,\%, 2,54\,MB en INT8 sur DPU
  \citep[Table~5]{2026-fata-low-cost-fpga-enhanced-cnn}), et le détecteur certifiable sur GPU
  de \citet{arxiv-2406.12362-certified-ml-object-detection-surveillance}.}}

\newglossaryentry{yolov4}{type=modeles, name={YOLOv4}, sort={YOLO 04}, parent=yolo,
  description={Quatrième version (2020) : backbone CSPDarknet53, module SPP et neck PANet,
  tête de YOLOv3 à ancres ; ajoute des techniques d'entraînement (Mosaic, CIoU) et
  l'activation Mish. Dans le corpus, Schlosser et al. la portent en INT8 sur DPU (ZCU104) et
  sur AGX Xavier : mAP COCO 37,5 contre 38,6, 1,31 contre 2,05\,FPS/W
  \citep[Table~1]{arxiv-2304.11580-framework-benchmarking-real-time-embedded} ; Feng 2022 la
  mesure sur GPU embarqués.}}

\newglossaryentry{yolov5}{type=modeles, name={YOLOv5}, sort={YOLO 05}, parent=yolo,
  description={Version d'Ultralytics (2020), sans article : réimplémentation en PyTorch, à
  ancres, backbone CSP, neck PANet, activation SiLU, déclinée en cinq tailles (n, s, m, l, x)
  par des multiplicateurs de profondeur et de largeur. Face à YOLOv4 : écosystème
  (entraînement, export ONNX) plus que rupture d'architecture. Dans le corpus : SATAY, Yan
  2024, et un YOLOv5 de taille non rapportée sur KV260 (Vitis AI).}}

\newglossaryentry{yolov5n}{type=modeles, name={YOLOv5n}, sort={YOLO 05 1n}, parent=yolov5,
  description={Taille \emph{nano}, la plus petite de YOLOv5. SATAY la déploie sur ZCU104 en
  640$\times$640 : 26,7 de mAP50 COCO, 21,41\,ms
  \citep[Table~III]{2023-montgomerie-corcoran-satay-streaming-architecture-toolflow-accelerating},
  seul résultat de SATAY sur SoC embarqué.}}

\newglossaryentry{yolov5s}{type=modeles, name={YOLOv5s}, sort={YOLO 05 2s}, parent=yolov5,
  description={Taille \emph{small} de YOLOv5. \citet{2024-yan-fpga-yolov5-accelerator-real-time}
  la quantifient entièrement en 4~bits (SiLU remplacée par OPReLU) sur ZCU102 : 79,1\,FPS à
  512$\times$512 pour 12,3\,W \citep[§7.4, Table~3]{2024-yan-fpga-yolov5-accelerator-real-time} ;
  SATAY la couvre aussi.}}

\newglossaryentry{yolov5m}{type=modeles, name={YOLOv5m}, sort={YOLO 05 3m}, parent=yolov5,
  description={Taille \emph{medium} de YOLOv5. N'apparaît dans le corpus que sur GPU de bureau,
  dans l'étude de prédiction conforme sur signaux ferroviaires
  \citep[§5.1.1]{arxiv-2304.06052-confident-object-detection-conformal-prediction}\pp, aux côtés
  de DETR.}}

\newglossaryentry{yolov8}{type=modeles, name={YOLOv8}, sort={YOLO 08}, parent=yolo,
  description={Version d'Ultralytics (2023) : tête découplée (classification et régression
  séparées) et \emph{sans ancres}, blocs C2f (avec une opération \emph{split}), activation
  SiLU, régression par Distribution Focal Loss ; tailles n à x. Face à YOLOv5 : abandon des
  ancres et tête découplée. Dans le corpus : SATAY, SMOF, Karki (sur DPU) et les mesures sur
  Jetson Orin de Rey et Alqahtani.}}

\newglossaryentry{yolov8n}{type=modeles, name={YOLOv8n}, sort={YOLO 08 1n}, parent=yolov8,
  description={Taille \emph{nano} de YOLOv8. Sur DPU, Karki et al. doivent remplacer SiLU par
  ReLU et le \emph{split} par des convolutions 1$\times$1 ; la mAP50-95 passe de 0,64 à 0,34 sur
  KITTI et de 0,60 à 0,45 sur VOC après PTQ INT8
  \citep[Table~II]{arxiv-2607.13106-no-attention-no-problem-dpu}\pp. SMOF la traite en bfp8 sur
  VCU118 ; Rey la mesure sur Orin Nano
  \citep[Table~4]{arxiv-2502.15737-performance-analysis-you-only-look}.}}

\newglossaryentry{yolov8s}{type=modeles, name={YOLOv8s}, sort={YOLO 08 2s}, parent=yolov8,
  description={Taille \emph{small} de YOLOv8. Plus gros modèle couvert par SATAY, mais sur
  VCU118 (W8A16) et non sur un SoC embarqué
  \citep[Table~III]{2023-montgomerie-corcoran-satay-streaming-architecture-toolflow-accelerating} ;
  aussi mesurée sur Jetson Orin
  \citep[Table~4]{arxiv-2502.15737-performance-analysis-you-only-look}.}}

\newglossaryentry{yolov11n}{type=modeles, name={YOLOv11n}, sort={YOLO 11n}, parent=yolo,
  description={Taille \emph{nano} de YOLO11 (Ultralytics, 2024), qui succède à YOLOv8 : blocs
  C3k2 et module d'attention C2PSA dans le backbone, tête sans ancres. L'attention et SiLU ne
  passent pas sur la DPU : Karki et al. approchent C2PSA par des convolutions 1$\times$1 et un
  hard-sigmoid \citep[§III-B]{arxiv-2607.13106-no-attention-no-problem-dpu}\pp. Le modèle perd
  2 points de mAP50-95 sur VOC (0,59 à 0,57) après INT8
  \citep[Table~II]{arxiv-2607.13106-no-attention-no-problem-dpu}\pp.}}

\newglossaryentry{yolov26n}{type=modeles, name={YOLOv26n}, sort={YOLO 26n}, parent=yolo,
  description={Taille \emph{nano} de YOLO26 (Ultralytics), présentée comme détecteur de bout
  en bout sans NMS et sans Distribution Focal Loss, pour faciliter l'export vers l'embarqué.
  Seul détecteur récent mesuré sur FPGA dans le corpus : sur DPU B4096 (ZCU104), après les
  mêmes adaptations que YOLOv11n, 34,05\,FPS de bout en bout sur une vidéo de personnes
  \citep[§IV-C]{arxiv-2607.13106-no-attention-no-problem-dpu}\pp ; le NMS y est conservé
  \citep[§III-B]{arxiv-2607.13106-no-attention-no-problem-dpu}\pp.}}

\newglossaryentry{gaussian-yolov3}{type=modeles, name={Gaussian YOLOv3}, sort={YOLO z Gaussian},
  parent=yolo,
  description={YOLOv3 dont chaque coordonnée de boîte est modélisée par une gaussienne
  (moyenne et variance) : la tête sort 8 valeurs au lieu de 4, la perte de boîte devient une
  log-vraisemblance négative, et l'incertitude de localisation pondère le score de détection
  \citep[§3.1--3.3]{arxiv-1904.04620-gaussian-yolov3-accurate-fast-object}\pp. Surcoût de
  0,04\,\% de FLOPs, +3,09 de mAP sur KITTI
  \citep[§3.1, Table~2]{arxiv-1904.04620-gaussian-yolov3-accurate-fast-object}\pp. Seule
  estimation d'incertitude de détecteur du corpus, mesurée sur GPU de bureau seulement.}}

\newglossaryentry{lp-yolo}{type=modeles, name={LP-YOLO}, sort={YOLO z LP}, parent=yolo,
  description={Détecteur YOLO allégé à une seule échelle (environ 0,35\,M paramètres,
  0,091\,GMAC), entraîné en QAT avec Brevitas (w4a4 ; w2a4 pour la variante Slim en
  256$\times$256) et compilé par FINN sur PYNQ-Z1
  \citep[§III-A, §IV-B]{arxiv-2609.24757-npu-accelerator-quantized-real-time}\pp. LP-YOLO Slim
  tient 35,66\,FPS pour 2,91\,W de carte avec 0,594 de mAP@0,5 VOC
  \citep[§V]{arxiv-2609.24757-npu-accelerator-quantized-real-time}\pp : seul détecteur FPGA du
  corpus mesuré sous 3\,W.}}

\newglossaryentry{req-yolo}{type=modeles, name={REQ-YOLO}, sort={YOLO z REQ}, parent=yolo,
  description={Cadre de compression de Tiny YOLO (YOLOv2-tiny) : matrices bloc-circulantes
  (convolution par FFT) et quantification hétérogène par ADMM, choisie couche par couche pour
  équilibrer DSP et LUT ; tous les poids tiennent en BRAM
  \citep[§3--5]{2019-ding-req-yolo-resource-aware-efficient}. Sur Virtex-7 XC7VX690T :
  314,2\,FPS pour 21\,W, latence de 3,18\,ms hors transferts et NMS
  \citep[Tables~1--2]{2019-ding-req-yolo-resource-aware-efficient}. Plateforme PCIe, non
  embarquée.}}

\newglossaryentry{tincy-yolo}{type=modeles, name={Tincy YOLO}, sort={YOLO z Tincy}, parent=yolo,
  description={Tiny YOLO modifié par FINN-R : poids binaires et activations 3~bits (W1A3) sauf
  première et dernière couches (8~bits), leaky-ReLU remplacé par ReLU ; 36\,\% d'opérations en
  moins pour une mAP VOC qui passe de 57,1 à 50,1
  \citep[§5.1]{arxiv-1809.04570-finn-r-end-end-deep}. Premier YOLO quantifié sur SoC Zynq du
  corpus, rapporté en GOp/s seulement (133 sur PYNQ-Z1, 288 sur Ultra96), sans FPS ni latence
  \citep[Table~5]{arxiv-1809.04570-finn-r-end-end-deep}.}}

\newglossaryentry{tinyissimoyolo}{type=modeles, name={TinyissimoYOLO}, sort={YOLO z Tinyissimo},
  parent=yolo,
  description={Détecteur pour microcontrôleur inspiré de YOLOv1 : entrée 88$\times$88,
  convolutions 3$\times$3 uniquement, une couche dense, sortie en grille 4$\times$4, INT8 par
  QAT, dimensionné pour les 442\,kB de poids du MAX78000
  \citep[§II-A]{arxiv-2306.00001-tinyissimoyolo-quantized-low-memory-footprint}. 58,5\,\% de mAP
  sur VOC réduit à 3 classes et 196\,µJ par inférence sur MAX78000, MCU seul
  \citep[§III, Table~II]{arxiv-2306.00001-tinyissimoyolo-quantized-low-memory-footprint}. Point
  de comparaison MCU des détecteurs FPGA.}}

% ---------------------------------------------------------- Autres détecteurs
\newglossaryentry{ssd}{type=modeles, name={SSD}, sort={SSD},
  description={\emph{Single Shot MultiBox Detector} : détecteur en une étape à ancres
  (\emph{default boxes}), qui prédit boîtes et classes sur plusieurs cartes de caractéristiques
  de résolutions décroissantes, puis applique un NMS. Dans le corpus : TGIIF, un SSD élagué en
  8~bits, gagne le volet FPGA de DAC-SDC'18 (IoU 0,62 à 11,96\,FPS sur XC7Z020)
  \citep[§5.2.1]{2021-xu-dac-sdc-low-power-object} ; SSD MobileNet V1 est mesuré sur cartes
  embarquées \citep[§IV]{arxiv-2409.16808-comprehensive-evaluation-deep-learning-object}\pp.}}

\newglossaryentry{centernet}{type=modeles, name={CenterNet}, sort={CenterNet},
  description={Détecteur \emph{sans ancres} (\enquote{Objects as Points}) : chaque objet est
  représenté par le pic de son centre sur une carte de chaleur par classe, la tête régressant
  taille et décalage ; les détections sont les maxima locaux, ce qui supprime le NMS. Dans le
  corpus, sa tête est reprise par CoDeNet, précisément pour éviter un NMS coûteux sur FPGA
  \citep[§4.1]{arxiv-2006.08357-codenet-efficient-deployment-input-adaptive}.}}

\newglossaryentry{codenet}{type=modeles, name={CoDeNet}, sort={CoDeNet}, parent=centernet,
  description={Détecteur co-conçu avec son accélérateur : backbone ShuffleNetV2, convolution
  déformable adaptée au matériel (depthwise, offsets bornés, carrés et entiers), tête CenterNet
  sans NMS, W4A8 \citep[§3.1, §4.1]{arxiv-2006.08357-codenet-efficient-deployment-input-adaptive}.
  Sur Ultra96 en 256$\times$256 : 51,1 d'AP50 VOC07 à 32,2\,FPS pour 5,6\,W de carte, modèle de
  0,76\,MB \citep[Table~3, §5]{arxiv-2006.08357-codenet-efficient-deployment-input-adaptive}.}}

\newglossaryentry{skynet}{type=modeles, name={SkyNet}, sort={SkyNet},
  description={Réseau de détection mono-objet conçu de bas en haut : un \enquote{Bundle}
  (DW-Conv3, PW-Conv1, BN, ReLU6) empilé 6 fois, choisi et dimensionné par recherche PSO
  contrainte par la latence GPU et FPGA, avec bypass et réordonnancement des cartes pour les
  petits objets, tête YOLO à 2 ancres
  \citep[§4--5]{arxiv-1909.09709-skynet-hardware-efficient-method-object}\pp. Vainqueur de
  DAC-SDC'19 sur TX2 et sur Ultra96 (IoU 0,716, 25,05\,FPS, 7,26\,W)
  \citep[Tables~5--6]{arxiv-1909.09709-skynet-hardware-efficient-method-object}\pp.}}

\newglossaryentry{detr}{type=modeles, name={DETR}, sort={DETR},
  description={\emph{DEtection TRansformer} : backbone CNN suivi d'un encodeur-décodeur
  transformer qui produit directement un ensemble fixe de boîtes à partir de requêtes apprises ;
  l'appariement hongrois de la perte supprime ancres et NMS. Dans le corpus, seulement sur GPU,
  comme l'un des trois détecteurs de l'étude de prédiction conforme ferroviaire
  \citep[§5.1.1]{arxiv-2304.06052-confident-object-detection-conformal-prediction}\pp. Ses
  opérations d'attention ne passent pas telles quelles sur une DPU.}}

% -------------------------------------------------------------- Backbones
\newglossaryentry{alexnet}{type=modeles, name={AlexNet}, sort={AlexNet},
  description={Classifieur ImageNet de 2012 : 5 couches de convolution (noyaux hétérogènes,
  jusqu'à 11$\times$11) puis 3 couches denses, ReLU. Banc d'essai historique des accélérateurs :
  \citet{2015-zhang-optimizing-fpga-accelerator-design-deep} atteignent 61,62\,GFLOPS sur ses
  convolutions en FP32 sur Virtex-7
  \citep[Tables~7--8]{2015-zhang-optimizing-fpga-accelerator-design-deep} ; EF-Train l'entraîne
  sur ZCU102.}}

\newglossaryentry{vgg16}{type=modeles, name={VGG16}, sort={VGG16},
  description={Classifieur de 2014 à 13 couches de convolution 3$\times$3 et 3 couches denses,
  très régulier mais lourd. Sa régularité convient aux moteurs de calcul uniques, à l'inverse
  d'AlexNet \citep[§2.7]{2018-venieris-toolflows-mapping-convolutional-neural-networks} ;
  EF-Train l'entraîne sur ZCU102 (46,99\,GFLOPS)
  \citep[Table~8]{arxiv-2202.10935-ef-train-enable-efficient-device}\pp.}}

\newglossaryentry{resnet}{type=modeles, name={ResNet}, sort={ResNet},
  description={Famille de réseaux à connexions résiduelles (ResNet-18, -50…), qui permettent
  d'entraîner des réseaux profonds ; les connexions longues imposent des tampons en dataflow.
  Dans le corpus : ResNet-50 W1A2 en dataflow FINN sur Alveo
  \citep[Table~IV]{arxiv-2011.07317-memory-efficient-dataflow-inference-deep} et Elastic-DF ;
  ResNet18 figé sur DPU comme extracteur de SLDA
  \citep[§III-A]{2021-piyasena-accelerating-continual-learning-edge-fpga} ; un ResNet-8 dérivé
  dans hls4ml.}}

\newglossaryentry{mobilenet}{type=modeles, name={MobileNet}, sort={MobileNet},
  description={Famille de backbones légers fondés sur la convolution séparable en profondeur
  (depthwise puis pointwise 1$\times$1) ; V2 ajoute les blocs résiduels inversés à goulot
  linéaire. Extracteur figé de plusieurs travaux d'apprentissage continu : MobileNetV1 dans
  Chameleon et QLR-CL, MobileNetV2 INT8 dans ImageHD
  \citep[§IV-A]{arxiv-2604.21280-imagehd-energy-efficient-device-continual} ; MobileNetV1 sert
  aussi de cas d'Elastic-DF.}}

\newglossaryentry{squeezenet}{type=modeles, name={SqueezeNet}, sort={SqueezeNet},
  description={Classifieur compact fondé sur des modules \enquote{fire} (compression 1$\times$1
  puis expansion 1$\times$1 et 3$\times$3). Dans le corpus, l'équipe SystemsETHZ de DAC-SDC le
  binarise et lui ajoute une tête YOLO : IoU de 0,49 à 25,97\,FPS sur XC7Z020
  \citep[§5.2.1]{2021-xu-dac-sdc-low-power-object}.}}

\newglossaryentry{mcunetv2}{type=modeles, name={MCUNetV2}, sort={MCUNetV2},
  description={Réseaux pour MCU obtenus par recherche conjointe d'architecture et
  d'ordonnancement : les premiers blocs sont exécutés par patchs pour abaisser le pic de SRAM,
  le champ récepteur est redistribué, déploiement INT8 avec TinyEngine
  \citep[§3]{arxiv-2110.15352-mcunetv2-memory-efficient-patch-inference}. Avec une tête YOLOv3,
  MCUNetV2-M4 atteint 64,6\,\% de mAP VOC en 247\,kB de SRAM
  \citep[Table~3]{arxiv-2110.15352-mcunetv2-memory-efficient-patch-inference}, sans latence ni
  énergie publiées.}}

% ------------------------------------------------------ Réseaux quantifiés
\newglossaryentry{qnn}{type=modeles, name={QNN}, sort={QNN},
  description={\emph{Quantized Neural Network} : réseau dont poids et activations sont en
  entiers de faible largeur (8~bits et moins), souvent entraîné en QAT. C'est la forme visée par
  FINN, FINN-R et hls4ml. Dans le corpus, la vérification formelle bit-exacte d'un QNN
  (encodage SMT) existe mais ne passe pas encore à l'échelle
  \citep{arxiv-2012.08185-scalable-verification-quantized-neural-networks}, et des fautes
  \enquote{stuck-at} sont injectées dans des QNN FINN sur ZCU104
  \citep{arxiv-1912.07394-efficient-error-tolerant-quantized-neural}.}}

\newglossaryentry{bnn}{type=modeles, name={BNN}, sort={BNN}, parent=qnn,
  description={\emph{Binarized Neural Network} : cas limite du QNN, poids et activations sur
  1~bit ($\pm 1$) ; le produit scalaire devient XNOR et comptage de bits (\emph{popcount}),
  la BN suivie du signe un simple seuil, ce qui convient aux LUT d'un FPGA. Base de FINN
  \citep{arxiv-1612.07119-finn-framework-fast-scalable-binarized} ; un BNN s'encode aussi
  exactement en SAT pour la vérification
  \citep[§4]{arxiv-1709.06662-verifying-properties-binarized-deep-neural}\pp.}}

\newglossaryentry{cnv}{type=modeles, name={CNV}, sort={CNV}, parent=bnn,
  description={Classifieur binaire de référence de FINN pour CIFAR-10 et SVHN : 3 fois (conv
  3$\times$3, conv 3$\times$3, max-pooling) avec 64, 128 et 256 canaux, puis 2 couches denses de
  512 \citep[§5.1]{arxiv-1612.07119-finn-framework-fast-scalable-binarized}. 80,1\,\% sur
  CIFAR-10 ; CNV-max tient 21\,900 images/s sur ZC706
  \citep[Tables~3--4]{arxiv-1612.07119-finn-framework-fast-scalable-binarized}. Réutilisé par
  FCMP et pour l'injection de fautes.}}

% ------------------------------------------------------------ Autoencodeurs
\newglossaryentry{ae}{type=modeles, name={AE (autoencodeur)}, sort={AE},
  description={Réseau entraîné à reconstruire son entrée au travers d'un goulot : l'encodeur
  la comprime en un code latent, le décodeur la reconstruit, avec une perte de reconstruction
  (MSE). Deux usages dans le corpus : la compression (seul l'encodeur est embarqué) et la
  détection d'anomalies (le score est l'erreur de reconstruction). Sur FPGA, il traite surtout
  des données non visuelles ou de petites images.}}

\newglossaryentry{cae}{type=modeles, name={CAE}, sort={CAE}, parent=ae,
  description={Autoencodeur convolutif : convolutions et sous-échantillonnage à l'encodeur,
  suréchantillonnage ou déconvolution au décodeur. Dans le corpus : détection d'artefacts
  d'images hyperspectrales par un CAE de 2,95\,M paramètres sur DPU (ZCU104)
  \citep[§III-B--D]{arxiv-2407.17647-energy-efficient-artefact-detection-accelerator}, CAE de
  vibrations de 8\,755 paramètres (Malviya), encodeurs de compression embarquée (Manjunath,
  Guerrisi), CAE débruiteur.}}

\newglossaryentry{vae}{type=modeles, name={VAE}, sort={VAE}, parent=ae,
  description={Autoencodeur variationnel : l'encodeur sort une moyenne $\mu$ et un écart-type
  $\sigma$, le latent est tiré par reparamétrisation, et la perte ajoute à la reconstruction une
  divergence de Kullback-Leibler vers une loi normale. Au déclenchement du LHC, seul l'encodeur
  est implanté et le score d'anomalie est la $D_{KL}$, en 80\,ns sur VU9P
  \citep{arxiv-2108.03986-autoencoders-fpgas-real-time-unsupervised} ; Que 2024 implante un VAE
  complet avec générateur gaussien.}}

\newglossaryentry{ssae}{type=modeles, name={SSAE}, sort={SSAE}, parent=ae,
  description={\emph{Stacked Sparse Autoencoder} : encodeurs d'autoencodeurs parcimonieux
  pré-entraînés sans supervision puis empilés, terminés par un softmax supervisé. Coutinho et
  al. implantent un SSAE 784-100-50-10 en réseau systolique virgule fixe sur Virtex-6 : 93,38\,\%
  sur MNIST contre 93,4\,\% sous Matlab
  \citep[§V-D]{2019-coutinho-deep-neural-network-hardware-implementation}.}}

\newglossaryentry{econ-ae}{type=modeles, name={ECON-AE}, sort={ECON-AE}, parent=ae,
  description={Autoencodeur de compression des données du calorimètre HGCAL de CMS (ECON-T) ;
  seul l'encodeur (convolution et couche dense) est embarqué, évalué par l'EMD. Campos et al.
  le quantifient en précision mixte guidée par la Hessienne puis le synthétisent avec hls4ml,
  sans mesure sur carte \citep[§7]{2024-campos-end-end-codesign-hessian-aware}.}}

\newglossaryentry{onlad}{type=modeles, name={ONLAD}, sort={ONLAD}, parent=ae,
  description={Détecteur d'anomalies en ligne : autoencodeur à une couche cachée entraîné par
  OS-ELM, avec un mécanisme d'oubli léger pour suivre la dérive, et implanté en virgule fixe
  32~bits sur XC7Z020 (ONLAD Core)
  \citep[§2--4]{2019-tsukada-neural-network-device-learning-anomaly}. Apprentissage continu en
  3,1\,W de carte \citep[§6.4]{2019-tsukada-neural-network-device-learning-anomaly}, sur des
  données tabulaires ou des images aplaties.}}

% ------------------------------------------------- Compression apprise d'images
\newglossaryentry{lic}{type=modeles, name={LIC}, sort={LIC},
  description={\emph{Learned Image Compression} : codec d'images appris, de type VAE avec
  hyper-prior (Ballé et al.) : transformée d'analyse vers un latent quantifié, hyper-analyse qui
  estime $\mu$ et $\sigma$ du latent pour le codage entropique, transformée de synthèse au
  décodage ; activations GDN. Sur SoC, \citet{arxiv-2503.04832-lightweight-embedded-fpga-deployment-learned}\pp{}
  distillent, élaguent et quantifient un LIC pour trois DPU B4096 (ZCU102) : 43,63\,FPS en
  encodage \citep[Table~VII]{arxiv-2503.04832-lightweight-embedded-fpga-deployment-learned}\pp.}}

\newglossaryentry{f-lic}{type=modeles, name={F-LIC}, sort={F-LIC}, parent=lic,
  description={Codec LIC à hyper-prior en virgule fixe 8~bits sur VCU118, en pipeline couche
  par couche ; deux poids partagent un DSP (casDSP) et la déconvolution saute les zéros ; le
  codage entropique reste sur l'hôte \citep{2022-sun-f-lic-fpga-learned-image}. Il encode du
  720p à 40,69\,FPS \citep[Fig.~6]{2022-sun-f-lic-fpga-learned-image} : premier codec appris
  temps réel sur FPGA du corpus, mais sur une grosse puce.}}

\newglossaryentry{gdn}{type=modeles, name={GDN}, sort={GDN}, parent=lic,
  description={\emph{Generalized Divisive Normalization} : activation des LIC qui divise chaque
  canal par la racine d'une combinaison pondérée des carrés des autres canaux ; l'inverse (iGDN)
  sert au décodeur. Non supportée par la DPU, elle est implantée comme opérateur XIR
  personnalisé en virgule fixe 32~bits
  \citep[§III-B, §III-D]{arxiv-2503.04832-lightweight-embedded-fpga-deployment-learned}\pp ; la
  remplacer par ReLU coûte 0,5 à 0,7\,dB de BD-PSNR
  \citep[Fig.~15]{arxiv-2503.04832-lightweight-embedded-fpga-deployment-learned}\pp.}}

% -------------------------------------------------- Détection d'anomalies (image)
\newglossaryentry{patchcore}{type=modeles, name={PatchCore}, sort={PatchCore},
  description={Détection d'anomalies industrielles par banque de caractéristiques : les
  descripteurs de patchs d'images normales, extraits par un CNN pré-entraîné, sont réduits en
  un \emph{coreset} ; le score est la distance au plus proche voisin. 99,1 d'AUROC image sur
  MVTec AD mais 92,2 sur un jeu réel d'huile alimentaire
  \citep[Tables~10, 12]{2023-cui-survey-unsupervised-anomaly-detection-algorithms}
  (Table~10 reconstruite par Docling) ; coûteux en mémoire, non porté sur FPGA dans le corpus.}}

\newglossaryentry{fastflow}{type=modeles, name={FastFlow}, sort={FastFlow},
  description={Détection d'anomalies par flot normalisant (\emph{normalizing flow}) 2D appris
  sur les caractéristiques d'un CNN figé : il estime la densité des images normales, et une
  faible vraisemblance signale l'anomalie. 99,4 d'AUROC image sur MVTec AD
  \citep[Table~10]{2023-cui-survey-unsupervised-anomaly-detection-algorithms} (reconstruite par
  Docling) ; comme PatchCore, absent des FPGA du corpus.}}

% ------------------------------------------------------- Apprentissage continu
\newglossaryentry{os-elm}{type=modeles, name={OS-ELM}, sort={OS-ELM},
  description={\emph{Online Sequential Extreme Learning Machine} : réseau à une couche cachée
  dont les poids d'entrée sont aléatoires et figés ; seuls les poids de sortie sont mis à jour
  en ligne par moindres carrés récursifs, sans rétropropagation. Avec un lot de 1, l'inversion
  matricielle devient une division scalaire
  \citep[§3.1]{2019-tsukada-neural-network-device-learning-anomaly}, ce qui rend l'apprentissage
  sur FPGA bon marché (ONLAD).}}

\newglossaryentry{slda}{type=modeles, name={SLDA}, sort={SLDA},
  description={\emph{Streaming Linear Discriminant Analysis} (Deep SLDA) : la dernière couche
  d'un CNN figé est remplacée par un classifieur linéaire qui met à jour en flux une moyenne
  par classe et une covariance partagée, sans rétropropagation ni replay
  \citep[§III-A]{2021-piyasena-accelerating-continual-learning-edge-fpga}. Piyasena et al.
  l'accélèrent en PL derrière un ResNet18 sur DPU : apprentissage à 168,67\,FPS contre 170,20
  en inférence seule \citep[Table~V]{2021-piyasena-accelerating-continual-learning-edge-fpga}
  (tableau reconstruit par Docling).}}

\newglossaryentry{chameleon}{type=modeles, name={Chameleon}, sort={Chameleon},
  description={Apprentissage continu en ligne par replay d'activations latentes (couche 21 de
  MobileNetV1) depuis deux tampons, court terme sur la puce et long terme hors puce ; seule la
  partie haute est rétropropagée, en FP16
  \citep[§III, §V-A]{2024-aggarwal-chameleon-dual-memory-replay-online}. Sur ZCU102 : 413\,ms et
  1,22\,J par image, 7 fois moins d'énergie que le Latent Replay, mais plus lent qu'une Jetson
  Nano \citep[Table~VI]{2024-aggarwal-chameleon-dual-memory-replay-online}.}}

\newglossaryentry{ef-train}{type=modeles, name={EF-Train}, sort={EF-Train},
  description={Accélérateur d'entraînement complet de CNN sur FPGA : un même noyau de
  convolution exécute propagation avant, rétropropagation et mise à jour, en FP32, avec une
  réorganisation des données en DRAM (\emph{data reshaping})
  \citep[§3--4]{arxiv-2202.10935-ef-train-enable-efficient-device}\pp. 2,08\,ms par image
  (batch 128) et 6,89\,W estimés sur ZCU102
  \citep[Table~7]{arxiv-2202.10935-ef-train-enable-efficient-device}\pp ; aucune stratégie
  contre l'oubli.}}

\newglossaryentry{qlr-cl}{type=modeles, name={QLR-CL}, sort={QLR-CL},
  description={\emph{Quantized Latent Replay} : un MobileNet-V1 est coupé en un étage figé
  quantifié en INT8 et un étage adaptatif réentraîné en FP32, nourri d'un mélange de nouvelles
  images et de \emph{latent replays} stockés quantifiés
  \citep[§III-C]{arxiv-2110.10486-tinyml-platform-device-continual-learning}. Sur l'ASIC PULP
  VEGA (22\,nm), 65 fois plus rapide et 37 fois plus économe qu'un STM32L4
  \citep[§V-E]{arxiv-2110.10486-tinyml-platform-device-continual-learning}.}}

\newglossaryentry{gdumb}{type=modeles, name={GDumb}, sort={GDumb},
  description={Méthode d'apprentissage continu \enquote{naïve} : un tampon équilibré par classe
  est rempli avidement, et le modèle est réentraîné sur ce seul tampon au moment de
  l'inférence. Il sert de référence de replay : 19,8\,\% sur CIFAR-100 en 10 tâches
  \citep[Table~S2]{arxiv-2010.15277-class-incremental-learning-survey-performance}\pp, et de
  politique de mémoire à TinyCL.}}

\newglossaryentry{tinycl}{type=modeles, name={TinyCL}, sort={TinyCL}, parent=gdumb,
  description={Accélérateur d'apprentissage continu en RTL : une unité de 9 blocs MAC
  réutilisée pour propagation avant, rétropropagation et mise à jour, en virgule fixe 16~bits,
  avec une mémoire de replay gérée selon GDumb
  \citep[§III-A, §III-C--F]{arxiv-2402.09780-tinycl-efficient-hardware-architecture-continual}\pp.
  Synthétisé en 65\,nm : 86\,mW d'après la synthèse, sans silicium
  \citep[Table~I]{arxiv-2402.09780-tinycl-efficient-hardware-architecture-continual}\pp.}}

\newglossaryentry{hdc}{type=modeles, name={HDC}, sort={HDC},
  description={\emph{Hyperdimensional Computing} : les données sont encodées en hypervecteurs
  de très grande dimension (souvent binaires), manipulés par liaison, superposition et
  permutation ; l'apprentissage se réduit à accumuler des prototypes et l'inférence à une
  recherche de similarité (distance de Hamming). Opérations de bits sans rétropropagation,
  adaptées au FPGA.}}

\newglossaryentry{imagehd}{type=modeles, name={ImageHD}, sort={ImageHD}, parent=hdc,
  description={Apprentissage continu non supervisé : un MobileNetV2 INT8 figé alimente un
  encodeur HDC binaire, les hypervecteurs sont regroupés en clusters dans une mémoire bornée,
  fusionnés par un kMeans++ simplifié ; accélérateur dataflow en Vitis HLS sur ZCU104
  \citep[§IV, §VI]{arxiv-2604.21280-imagehd-energy-efficient-device-continual}. 0,27\,ms et
  0,97\,mJ par image sur CIFAR-10, puissance estimée
  \citep[Table~VII]{arxiv-2604.21280-imagehd-energy-efficient-device-continual}.}}

\newglossaryentry{lwf}{type=modeles, name={LwF}, sort={LwF},
  description={\emph{Learning without Forgetting} : régularisation par distillation, sans
  exemplaires ; pendant l'apprentissage d'une nouvelle tâche, les sorties de l'ancien modèle
  sur les nouvelles données servent de cibles douces. Meilleure méthode sans exemplaires du
  banc de \citet{arxiv-2010.15277-class-incremental-learning-survey-performance}\pp{} (30,2\,\%
  sur CIFAR-100 en 10 tâches, Table~2), mais en dessous du simple fine-tuning avec 2\,000
  exemplaires (37,9\,\%) ; baseline de Chameleon.}}

\newglossaryentry{rilod}{type=modeles, name={RILOD}, sort={RILOD}, parent=lwf,
  description={Apprentissage incrémental pour détecteur en une étape (RetinaNet) : une nouvelle
  classe est ajoutée avec ses seules données, en distillant depuis l'ancien modèle figé les
  sorties de classification, les boîtes et des caractéristiques intermédiaires
  \citep[§4.2]{arxiv-1904.00781-rilod-near-real-time-incremental}\pp. Exécuté sur Jetson TX2
  ou en mode edge-cloud, sans FPGA : seul travail de détection incrémentale du corpus.}}

% ------------------------------------------------------ Opérateurs
\newglossaryentry{nms}{type=modeles, name={NMS}, sort={NMS},
  description={\emph{Non-Maximum Suppression} : post-traitement des détecteurs à ancres ou à
  grille, qui trie les boîtes par score et supprime celles dont l'IoU avec une boîte mieux
  notée dépasse un seuil. Séquentiel et dépendant des données, il tourne presque toujours sur
  le cœur ARM, hors des chiffres de l'accélérateur : REQ-YOLO l'exclut de sa latence
  \citep[Table~1]{2019-ding-req-yolo-resource-aware-efficient}, et il devient le goulot de
  bout en bout chez \citet{arxiv-2607.13106-no-attention-no-problem-dpu}\pp{} (§IV-C).}}

\newglossaryentry{relu}{type=modeles, name={ReLU}, sort={ReLU},
  description={\emph{Rectified Linear Unit}, $\max(0,x)$ ; variantes ReLU6 (écrêtée à 6) et
  leaky-ReLU (pente faible pour $x<0$, utilisée par YOLOv2 à v4). Linéaire par morceaux, elle
  coûte une comparaison et se fusionne avec la quantification ; c'est l'activation de repli des
  toolflows : FINN-R remplace le leaky-ReLU de Tincy YOLO par ReLU
  \citep[§5.1]{arxiv-1809.04570-finn-r-end-end-deep}, Karki et al. y ramènent SiLU pour la DPU.}}

\newglossaryentry{silu}{type=modeles, name={SiLU}, sort={SiLU},
  description={\emph{Sigmoid Linear Unit} (ou Swish), $x\cdot\sigma(x)$ : activation lisse de
  YOLOv5, v8 et suivants. Elle exige une sigmoïde, absente des opérateurs natifs de la DPU et
  de FINN \citep[§VI-D]{arxiv-2206.07527-qonnx-representing-arbitrary-precision-quantized} ; d'où
  des remplacements : ReLU chez Karki et al.
  \citep[§III-B]{arxiv-2607.13106-no-attention-no-problem-dpu}\pp, HardSwish dans SATAY
  \citep[§III-B]{2023-montgomerie-corcoran-satay-streaming-architecture-toolflow-accelerating},
  OPReLU chez Yan et al. \citep[§4]{2024-yan-fpga-yolov5-accelerator-real-time}.}}
